\subsection{有限の途中まで積分してから極限をとる}

$\int_0^1\frac1{\sqrt{x}}\dd x$では，$x=0$で被積分関数が定義されず，その近くで値が限りなく大きくなる。それでも，$\varepsilon>0$で切り取った区間では
\[
 \int_{\varepsilon}^1\frac1{\sqrt{x}}\dd x
 =2-2\sqrt{\varepsilon}
\]
と計算できる。この値が$\varepsilon\to0+$で$2$へ近づくことを利用して，積分の意味を広げる。

\begin{definition}{広義積分}
$f$が$[a,b)$の各閉部分区間で積分可能で，端点$b$に特異性をもつとき，
\[
 \int_a^bf(x)\dd x=\lim_{c\to b-}\int_a^cf(x)\dd x
\]
と定める。ただし，右辺が有限の実数として存在するとき，この広義積分は収束するという。左端に特異性がある場合も同様に，一側極限で定める。

上端が無限大の場合は
\[
 \int_a^{\infty}f(x)\dd x=\lim_{R\to\infty}\int_a^Rf(x)\dd x
\]
と定める。下端が$-\infty$の場合も同様である。有限の極限が存在しなければ発散するという。
\end{definition}

ここで，$+\infty$へ発散する場合も，振動して極限がない場合も，「収束しない」という点では同じである。一方，発散の仕方を区別したいときには「$+\infty$へ発散する」などと具体的に述べる。

